% Options for packages loaded elsewhere
\PassOptionsToPackage{unicode}{hyperref}
\PassOptionsToPackage{hyphens}{url}
\PassOptionsToPackage{dvipsnames,svgnames,x11names}{xcolor}
%
\documentclass[
  letterpaper,
  DIV=11,
  numbers=noendperiod]{scrreprt}

\usepackage{amsmath,amssymb}
\usepackage{iftex}
\ifPDFTeX
  \usepackage[T1]{fontenc}
  \usepackage[utf8]{inputenc}
  \usepackage{textcomp} % provide euro and other symbols
\else % if luatex or xetex
  \usepackage{unicode-math}
  \defaultfontfeatures{Scale=MatchLowercase}
  \defaultfontfeatures[\rmfamily]{Ligatures=TeX,Scale=1}
\fi
\usepackage{lmodern}
\ifPDFTeX\else  
    % xetex/luatex font selection
\fi
% Use upquote if available, for straight quotes in verbatim environments
\IfFileExists{upquote.sty}{\usepackage{upquote}}{}
\IfFileExists{microtype.sty}{% use microtype if available
  \usepackage[]{microtype}
  \UseMicrotypeSet[protrusion]{basicmath} % disable protrusion for tt fonts
}{}
\makeatletter
\@ifundefined{KOMAClassName}{% if non-KOMA class
  \IfFileExists{parskip.sty}{%
    \usepackage{parskip}
  }{% else
    \setlength{\parindent}{0pt}
    \setlength{\parskip}{6pt plus 2pt minus 1pt}}
}{% if KOMA class
  \KOMAoptions{parskip=half}}
\makeatother
\usepackage{xcolor}
\setlength{\emergencystretch}{3em} % prevent overfull lines
\setcounter{secnumdepth}{5}
% Make \paragraph and \subparagraph free-standing
\makeatletter
\ifx\paragraph\undefined\else
  \let\oldparagraph\paragraph
  \renewcommand{\paragraph}{
    \@ifstar
      \xxxParagraphStar
      \xxxParagraphNoStar
  }
  \newcommand{\xxxParagraphStar}[1]{\oldparagraph*{#1}\mbox{}}
  \newcommand{\xxxParagraphNoStar}[1]{\oldparagraph{#1}\mbox{}}
\fi
\ifx\subparagraph\undefined\else
  \let\oldsubparagraph\subparagraph
  \renewcommand{\subparagraph}{
    \@ifstar
      \xxxSubParagraphStar
      \xxxSubParagraphNoStar
  }
  \newcommand{\xxxSubParagraphStar}[1]{\oldsubparagraph*{#1}\mbox{}}
  \newcommand{\xxxSubParagraphNoStar}[1]{\oldsubparagraph{#1}\mbox{}}
\fi
\makeatother


\providecommand{\tightlist}{%
  \setlength{\itemsep}{0pt}\setlength{\parskip}{0pt}}\usepackage{longtable,booktabs,array}
\usepackage{calc} % for calculating minipage widths
% Correct order of tables after \paragraph or \subparagraph
\usepackage{etoolbox}
\makeatletter
\patchcmd\longtable{\par}{\if@noskipsec\mbox{}\fi\par}{}{}
\makeatother
% Allow footnotes in longtable head/foot
\IfFileExists{footnotehyper.sty}{\usepackage{footnotehyper}}{\usepackage{footnote}}
\makesavenoteenv{longtable}
\usepackage{graphicx}
\makeatletter
\newsavebox\pandoc@box
\newcommand*\pandocbounded[1]{% scales image to fit in text height/width
  \sbox\pandoc@box{#1}%
  \Gscale@div\@tempa{\textheight}{\dimexpr\ht\pandoc@box+\dp\pandoc@box\relax}%
  \Gscale@div\@tempb{\linewidth}{\wd\pandoc@box}%
  \ifdim\@tempb\p@<\@tempa\p@\let\@tempa\@tempb\fi% select the smaller of both
  \ifdim\@tempa\p@<\p@\scalebox{\@tempa}{\usebox\pandoc@box}%
  \else\usebox{\pandoc@box}%
  \fi%
}
% Set default figure placement to htbp
\def\fps@figure{htbp}
\makeatother
% definitions for citeproc citations
\NewDocumentCommand\citeproctext{}{}
\NewDocumentCommand\citeproc{mm}{%
  \begingroup\def\citeproctext{#2}\cite{#1}\endgroup}
\makeatletter
 % allow citations to break across lines
 \let\@cite@ofmt\@firstofone
 % avoid brackets around text for \cite:
 \def\@biblabel#1{}
 \def\@cite#1#2{{#1\if@tempswa , #2\fi}}
\makeatother
\newlength{\cslhangindent}
\setlength{\cslhangindent}{1.5em}
\newlength{\csllabelwidth}
\setlength{\csllabelwidth}{3em}
\newenvironment{CSLReferences}[2] % #1 hanging-indent, #2 entry-spacing
 {\begin{list}{}{%
  \setlength{\itemindent}{0pt}
  \setlength{\leftmargin}{0pt}
  \setlength{\parsep}{0pt}
  % turn on hanging indent if param 1 is 1
  \ifodd #1
   \setlength{\leftmargin}{\cslhangindent}
   \setlength{\itemindent}{-1\cslhangindent}
  \fi
  % set entry spacing
  \setlength{\itemsep}{#2\baselineskip}}}
 {\end{list}}
\usepackage{calc}
\newcommand{\CSLBlock}[1]{\hfill\break\parbox[t]{\linewidth}{\strut\ignorespaces#1\strut}}
\newcommand{\CSLLeftMargin}[1]{\parbox[t]{\csllabelwidth}{\strut#1\strut}}
\newcommand{\CSLRightInline}[1]{\parbox[t]{\linewidth - \csllabelwidth}{\strut#1\strut}}
\newcommand{\CSLIndent}[1]{\hspace{\cslhangindent}#1}

\usepackage{booktabs}
\usepackage{caption}
\usepackage{longtable}
\usepackage{colortbl}
\usepackage{array}
\usepackage{anyfontsize}
\usepackage{multirow}
\KOMAoption{captions}{tableheading}
\makeatletter
\@ifpackageloaded{tcolorbox}{}{\usepackage[skins,breakable]{tcolorbox}}
\@ifpackageloaded{fontawesome5}{}{\usepackage{fontawesome5}}
\definecolor{quarto-callout-color}{HTML}{909090}
\definecolor{quarto-callout-note-color}{HTML}{0758E5}
\definecolor{quarto-callout-important-color}{HTML}{CC1914}
\definecolor{quarto-callout-warning-color}{HTML}{EB9113}
\definecolor{quarto-callout-tip-color}{HTML}{00A047}
\definecolor{quarto-callout-caution-color}{HTML}{FC5300}
\definecolor{quarto-callout-color-frame}{HTML}{acacac}
\definecolor{quarto-callout-note-color-frame}{HTML}{4582ec}
\definecolor{quarto-callout-important-color-frame}{HTML}{d9534f}
\definecolor{quarto-callout-warning-color-frame}{HTML}{f0ad4e}
\definecolor{quarto-callout-tip-color-frame}{HTML}{02b875}
\definecolor{quarto-callout-caution-color-frame}{HTML}{fd7e14}
\makeatother
\makeatletter
\@ifpackageloaded{bookmark}{}{\usepackage{bookmark}}
\makeatother
\makeatletter
\@ifpackageloaded{caption}{}{\usepackage{caption}}
\AtBeginDocument{%
\ifdefined\contentsname
  \renewcommand*\contentsname{Table of contents}
\else
  \newcommand\contentsname{Table of contents}
\fi
\ifdefined\listfigurename
  \renewcommand*\listfigurename{List of Figures}
\else
  \newcommand\listfigurename{List of Figures}
\fi
\ifdefined\listtablename
  \renewcommand*\listtablename{List of Tables}
\else
  \newcommand\listtablename{List of Tables}
\fi
\ifdefined\figurename
  \renewcommand*\figurename{Figure}
\else
  \newcommand\figurename{Figure}
\fi
\ifdefined\tablename
  \renewcommand*\tablename{Table}
\else
  \newcommand\tablename{Table}
\fi
}
\@ifpackageloaded{float}{}{\usepackage{float}}
\floatstyle{ruled}
\@ifundefined{c@chapter}{\newfloat{codelisting}{h}{lop}}{\newfloat{codelisting}{h}{lop}[chapter]}
\floatname{codelisting}{Listing}
\newcommand*\listoflistings{\listof{codelisting}{List of Listings}}
\makeatother
\makeatletter
\makeatother
\makeatletter
\@ifpackageloaded{caption}{}{\usepackage{caption}}
\@ifpackageloaded{subcaption}{}{\usepackage{subcaption}}
\makeatother

\usepackage{bookmark}

\IfFileExists{xurl.sty}{\usepackage{xurl}}{} % add URL line breaks if available
\urlstyle{same} % disable monospaced font for URLs
\hypersetup{
  pdftitle={Evaluation of Sharknado!},
  pdfauthor={Colton Haight},
  colorlinks=true,
  linkcolor={blue},
  filecolor={Maroon},
  citecolor={Blue},
  urlcolor={Blue},
  pdfcreator={LaTeX via pandoc}}


\title{Evaluation of Sharknado!}
\author{Colton Haight}
\date{}

\begin{document}
\maketitle

\renewcommand*\contentsname{Table of contents}
{
\hypersetup{linkcolor=}
\setcounter{tocdepth}{2}
\tableofcontents
}

\bookmarksetup{startatroot}

\chapter{Executive Summary}\label{executive-summary}

\section{Purpose of Report}\label{purpose-of-report}

This study dives into the possibility of a Sharknado event. A Sharknado
can be explained as a ``shark-infested tornado''. This is a tornado that
generates enough wind speed to lift a great white shark from the water
and into its vortex. The purpose of this study is to assess whether a
Sharknado event could be considered a risk for property damage insurance
payouts.

\section{Methods}\label{methods}

To start, a comparison of the weight of a great white shark (around 4500
lbs) and vehicles (around 4000 - 4600 lbs for SUVs and Trucks) is used
because cars are known to be lifted and destroyed by tornadoes. Using
the EF (Enhanced Fujita) Scale, the strength of tornadoes able to lift
and toss cars (or a great white shark) is found to be between EF3 - EF5
tornadoes. These are wind speeds of 136 mph to 200+ mph. This study will
only be looking at tornadoes of EF3 - EF5 for this reason. This study
also looks at four locations of tornado data from 2015 to 2026. These
locations are Florida, South Carolina, Virginia, Oklahoma, California,
and Los Angeles County, CA. These places were chosen because of their
abundant tropical storms and coastal area (Florida and South Carolina),
historic tornado amounts and strength (Oklahoma), and the setting of the
movie Sharknado (Los Angeles, CA).

\section{Findings}\label{findings}

Tornadoes of strength EF3 - EF5 are fairly uncommon for coastal states
like Florida and South Carolina. Florida had only five tornadoes in the
last 10 years that were of strength EF3, and zero EF4 and EF5 tornadoes.
All but one of these tornadoes started and ended on land. The one
tornado started on a beach in Florida, where the water is not deep
enough for sharks. South Carolina had nine EF3 and one EF4 tornado over
the last 10 years. All of these tornadoes started and stayed on land,
making a Sharknado impossible. Virginia had five EF3 tornadoes that all
happened on land. Oklahoma has had plenty of tornadoes strong enough to
lift and carry a shark, but there are no sharks in Oklahoma besides
those in the Oklahoma Aquarium in Jenks, OK. California has not had a
tornado above an EF2 over the last 10 years.

\section{Conclusions}\label{conclusions}

The probability of a Sharknado event is extremely low, if not a zero
percent chance. There was one tornado in the three coastal states that
might have had the ability to lift a shark. This South Carolina tornado
happened to start and end on land, making a Sharknado impossible. The
Oklahoma Aquarium houses sharks in a state that is known for its number
and strength of tornadoes. The Oklahoma Aquarium is in Tulsa County, OK,
where only EF2 or lower tornadoes have touched down in the last 10
years, rendering this very improbable. From the findings of this study,
the probability of a Sharknado happening is almost, if not, a zero
percent chance.

\section{Recommendations}\label{recommendations}

From this study, Sharknados are certainly extremely unlikely, if not a
zero percent chance. Even though strong storms and tornadoes without
sharks are still a cause for concern. Property damage and its resulting
insurance payouts will still be present with storms and tornadoes that
have high wind speeds and extreme precipitation events. These are real
and worrisome outcomes that happen all over the country that are more
probable than a Sharknado and should be taken more into account.

\bookmarksetup{startatroot}

\chapter{Introduction}\label{introduction}

This reports evaluates the potential insurance risk associated with a
Sharknado scenario for property damage payouts . A Sharknado is
explained by a tornado that generates enough wind speed to be able to
lift a heavy animal, such as a great white shark, out of the water and
into its vortex. A Sharknado can also be called a ``Shark-infested
tornado''.

This study starts by comparing the average great white shark to objects
of similar weight, specifically cars or pickup trucks. Tornadoes are
known for destroying homes and tossing around vehicles, thus finding the
types of vehicles that are similar enough in weight to a great white
shark makes finding the right magnitude of tornadoes required to lift a
great white shark easier.

Next, finding the magnitude of a tornado capable of picking up a great
white shark, or a car, is found using the EF scale (Enhanced Fujita
Scale) by looking at EF ratings (US Department of Commerce (n.d.)). The
Enhanced Fujita Scale uses ratings of EF0 to EF5. Tornadoes of interest
for this study are tornadoes of rating EF3 to EF5. These tornadoes have
strong enough wind speeds to potentially pick up and throw cars or, in
this case, a great white shark.

Finally, looking at places that are near the coast to be able to pick up
a great white shark is essential to actually generating a Sharknado.
There are 4 locations observed for their tornado data that fit this
agenda in this study: Florida, South Carolina, Virginia, California, and
Los Angeles (this is the setting of the Sharknado movie). For another
comparison of tornado strength and amount, Oklahoma is also looked at
for its tornado EF rating, even though it is landlocked and does not
have any native shark population. These locations have tornado data
taken from 2015 - 2026 using NOAA tornado data.

Combining the magnitude of tornadoes with the areas where these
tornadoes touch down gives an idea of the possibility of a potential
Sharknado event occurring.

\bookmarksetup{startatroot}

\chapter{Data Documentation}\label{data-documentation}

All data was collected from NOAA (National Oceanic and Atmospheric
Administration) Storm Events Database. The NOAA's storm events database
contains records on various types fo severe weather collected by NOAA's
National Weather Service.

Tornado data from NOAA Storm Events Database was collected for five
different states and one county across the United States ({``Storm
Events Database {\textbar} National Centers for Environmental
Information''} (n.d.)). These locations were chosen because of their
amount of EF3 or higher tornadoes, location close to the ocean, and
setting of the Sharknado movie. These locations are:

\begin{itemize}
\item
  Florida: high frequency of tornadoes EF3 or higher for a coastal
  region
\item
  South Carolina: high frequency of tornadoes EF3 or higher for a
  coastal region
\item
  Virginia: high frequency of tornadoes EF3 or higher for a coastal
  region
\item
  Oklahoma: high amount of powerful tornadoes to compare the frequency
  of high magnitude tornadoes
\item
  California and Los Angeles, CA: State and city setting of the movie
  ``Sharknado'' ({``Sharknado ({TV} Movie 2013) - Filming \& Production
  - {IMDb}''} (2013))
\end{itemize}

All data was analyzed using the programming software R and these
following R packages:

\begin{itemize}
\item
  \texttt{tidyverse}: data cleaning, data manipulation, and table
  formulation
\item
  \texttt{gt}: to make better looking tables
\item
  \texttt{maps}: make maps of United States for tornado data
\end{itemize}

\bookmarksetup{startatroot}

\chapter{Methods}\label{methods-1}

A Sharknado event involves a tornado being able to lift a great white
shark out of the ocean and carry it onto land. To see if this is
possible, we need to look at the physics: the weight of a great white
shark, what objects are known to be picked up by tornadoes, and the
magnitude of tornadoes required to pick up a shark. First, finding out
how a great white shark could be picked up by a tornado is nessesary.
The weight of a great white shark will be compared to other objects that
can and have been picked up by tornadoes, such as cars. Next, figuring
out what strength (in EF scale from EF3 - EF5) of a tornado would be
needed to be able to lift a great white shark and a car. Lastly,
identifying tornadoes from 2015-2026 that are in that EF range from 6
different places in the United States. These places are be Florida,
South Carolina, Virginia, Oklahoma, California, and Los Angeles. These
places were chosen for specific reasons. Florida has many tropical
storms and is a coastal state (Florida, (n.d.)). South Carolina is also
another place that can be hit by tropical storms and is another coastal
state (South Carolina, (n.d.)). Virginia is a coastal state that has had
a good amount of tornadoes of our criteria (Virginia, (n.d.b)). Oklahoma
is not a coastal state, but it is the tornado capital of the United
States, so we will just look at the number of tornadoes that meet the
criteria to get an idea of how rare strong tornadoes are (Oklahoma,
(n.d.)). California ((n.d.a)) was chosen because it is the state setting
of Sharknado, while Los Angeles ((n.d.)) was picked because it is the
city setting of the Sharknado movie ({``Sharknado ({TV} Movie 2013) -
Filming \& Production - {IMDb}''} (2013)).

Sharks and cars weigh relatively the same in pounds (lbs). The average
great white shark weighs around 4,500 lbs (Fisheries (2025)). A car SUV
weighs around 3,986 lbs while a truck SUV weighs around 4,578 lbs
(AutoInsurance (n.d.)). This means that your average great white shark
weighs more than a big car and as much as a truck or SUV. We will use a
car as the base to find what EF scale level will be used to lift a great
white shark out of the water.

The EF scale is made up of 6 different ratings of tornado based on
3-second wind speed in mph (US Department of Commerce (n.d.)). An EF0
tornado has the lowest wind speed, ranging between 65-85 mph, EF1 has
wind speeds between 86-110 mph, EF2 has wind speeds between 111-135 mph,
EF3 has wind speeds of 136-165 mph, EF4 has speeds between 166-200 mph,
and EF5 has wind speeds that are over 200 mph. An EF3 tornado is likely
to toss cars and trucks, while an EF4 is very likely to toss cars many
feet, and an EF5 tornado is almost certain to lift and destroy cars. EF3
tornadoes are unlikely to fully pick up a shark but it is still
important to analyze them because a great white shark could be moved by
an EF3 tornado. Because of the comparison of a great white shark to the
size and weight of a car SUV, and a truck SUV, that gives a good
comparison to be able to put a shark in the place of such cars in the EF
scales' ability to lift cars (Quach (2025)). Tornadoes from EF3 to EF5
are chosen because this is the strength of tornadoes that are able to
lift a car (Salaita (n.d.)).

Using these metrics, finding storms that have had the potential to
produce a Sharknado allows the ability to evaluate the potential
insurance risk with such events.

\bookmarksetup{startatroot}

\chapter{Results}\label{results}

The results of this study examines tornado activity over the last 10
years from specific states across the United States. Tornado strength of
EF3 - EF5 is chosen due to their capabilities to lift and toss objects
that are similar to sharks weight wise, these objects being cars or
trucks. The findings are shown by state and area (FL, SC, VA, CA, Los
Angeles) to be able to see where and which tornadoes had the strength
and location to start a Sharknado event. California and Los Angeles, CA
were chosen due to being the location of the film ``Sharknado''.

Below are these tornadoes being mapped on the entire United States to
get an idea of where these states and their tornadoes are and why they
were chosen for this study.

\begin{figure}

\centering{

\pandocbounded{\includegraphics[keepaspectratio]{results_files/figure-pdf/fig-plot1-1.pdf}}

}

\caption{\label{fig-plot1}EF3 - EF5 Tornado Data Accross USA}

\end{figure}%

Most tornadoes happen on the eastern side of the country as shown in
Figure~\ref{fig-plot1}. This shows the regions to check to see if
Sharknadoes are possible. The east coast is the best place to start the
search. The states that have enough tornadoes, and the states that will
be examined, are Florida, South Carolina, and Virginia.

\section{Results: Florida}\label{results-florida}

From Figure~\ref{fig-plot1}, Florida is a place that is on the coast.
This would make it possible to pick up a shark in the event of a strong
enough tornado. There are enough tornadoes that touched down in Florida
over the last 10 years which is a reason Florida is a state that will be
analyzed.

\begin{figure}

\centering{

\begin{table}
\fontsize{12.0pt}{14.0pt}\selectfont
\begin{tabular*}{\linewidth}{@{\extracolsep{\fill}}llrrr}
\toprule
STATE\_ABBR & EVENT\_TYPE & TOR\_F\_SCALE & BEGIN\_LAT & BEGIN\_LON \\ 
\midrule\addlinespace[2.5pt]
FL & Tornado & 3 & 30.8707 & -87.3810 \\ 
FL & Tornado & 3 & 30.4907 & -87.2052 \\ 
FL & Tornado & 3 & 30.4934 & -84.1262 \\ 
FL & Tornado & 3 & 30.6376 & -85.5108 \\ 
FL & Tornado & 3 & 30.1384 & -85.7526 \\ 
FL & Tornado & 3 & 26.8878 & -81.0381 \\ 
FL & Tornado & 3 & 26.5555 & -80.3321 \\ 
FL & Tornado & 3 & 27.3740 & -80.4270 \\ 
\bottomrule
\end{tabular*}
\end{table}

}

\caption{\label{fig-table1}Florida EF3 - EF5 Tornado Data}

\end{figure}%

Figure~\ref{fig-table1} shows that there have been only 5 EF3 tornadoes
in the last 10 years in the state of Florida. Because of an EF3 tornado
being the highest EF rating, and it being ``likely to toss a car or
truck'', a tornado in Florida is fairly unlikely to pick up the average
great white shark. Having only five EF3 tornadoes in the past 10 years
seems like the chances of a Sharknado happening seem pretty slim. Below
is a map showing these tornadoes.

\begin{figure}

\centering{

\pandocbounded{\includegraphics[keepaspectratio]{results_files/figure-pdf/fig-plot2-1.pdf}}

}

\caption{\label{fig-plot2}Florida EF3 - EF5 Tornado Map}

\end{figure}%

Figure~\ref{fig-plot2} shows a map of Florida showing the five EF3
tornadoes that have occurred in the last 10 years. There is one tornado
that shows up close to the water, but looking at the coordinates, it
appears on the beach. EF3 tornadoes are likely to toss cars and trucks
(Quach (2025)). Tossing and picking up two different things in this
scenario, leading to believe an EF3 tornado is not very likely to fully
pick up a great white shark. The EF3 tornadoes can displace the shark if
its close enough to land but that is still fairly unlikely.

\section{Results: South Carolina}\label{results-south-carolina}

A place with similar properties to Florida, having tropical storms
frequently and being a coastal state, is South Carolina. The data below
gives another state to compare to that is in an area where great white
sharks may be present.

\begin{figure}

\centering{

\begin{table}
\fontsize{12.0pt}{14.0pt}\selectfont
\begin{tabular*}{\linewidth}{@{\extracolsep{\fill}}llrrr}
\toprule
STATE\_ABBR & EVENT\_TYPE & TOR\_F\_SCALE & BEGIN\_LAT & BEGIN\_LON \\ 
\midrule\addlinespace[2.5pt]
SC & Tornado & 3 & 34.6170 & -83.0840 \\ 
SC & Tornado & 3 & 33.2751 & -81.7048 \\ 
SC & Tornado & 3 & 33.2754 & -81.5438 \\ 
SC & Tornado & 3 & 33.3177 & -81.2903 \\ 
SC & Tornado & 3 & 33.2071 & -81.2742 \\ 
SC & Tornado & 3 & 33.4539 & -81.2651 \\ 
SC & Tornado & 4 & 32.7045 & -81.2899 \\ 
SC & Tornado & 3 & 33.1816 & -79.9901 \\ 
SC & Tornado & 3 & 32.9286 & -81.3756 \\ 
SC & Tornado & 3 & 33.1154 & -81.1889 \\ 
\bottomrule
\end{tabular*}
\end{table}

}

\caption{\label{fig-table2}South Carolina EF3 - EF5 Tornado Data}

\end{figure}%

South Carolina is similar to Florida from Figure~\ref{fig-table2}. There
are a lot of EF3 tornadoes with one EF4 tornado. But it's in the same
boat as Florida; there has been a very limited number of powerful
tornadoes over the past 10 years in South Carolina. There have been nine
EF3 tornadoes and one EF4 tornado over that span. The surprising thing
is that the EF4 and six of the EF3 tornadoes all happened on the same
day. But the possibility of a Sharknado is very low because these
tornadoes all happened away from the coast. The EF4 happened in Estill
and Hampton Counties (Climatology Office (2020)). Below is a map showing
all the tornadoes from South Carolina over the last 10 years.

\begin{figure}

\centering{

\pandocbounded{\includegraphics[keepaspectratio]{results_files/figure-pdf/fig-plot3-1.pdf}}

}

\caption{\label{fig-plot3}South Carolina EF3 - EF5 Tornado Map}

\end{figure}%

The EF4 tornado is the main Sharknado threat. As seen in
Figure~\ref{fig-plot3}, all tornadoes EF3 and above only appear on land
and are far enough from the coast to have any real threat of a
Sharknado. From the evidence from Florida and now South Carolina, a
Sharknado is looking less and less possible. The insurance payout risk
associated with a Sharknado in South Carolina is negligible.

\section{Results: Virginia}\label{results-virginia}

Another coastal state similar to South Carolina is Virginia.
Figure~\ref{fig-plot1} shows a plethora of tornadoes in VA over the last
10 years that fit the criteria of EF3 or higher tornadoes to test the
potential of a Sharknado. Below is Virginia tornado data from the last
10 years.

\begin{figure}

\centering{

\begin{table}
\fontsize{12.0pt}{14.0pt}\selectfont
\begin{tabular*}{\linewidth}{@{\extracolsep{\fill}}llrrr}
\toprule
STATE\_ABBR & EVENT\_TYPE & TOR\_F\_SCALE & BEGIN\_LAT & BEGIN\_LON \\ 
\midrule\addlinespace[2.5pt]
VA & Tornado & 3 & 37.2318 & -78.8576 \\ 
VA & Tornado & 3 & 37.8285 & -76.9681 \\ 
VA & Tornado & 3 & 37.4586 & -79.1838 \\ 
VA & Tornado & 3 & 36.8564 & -79.8906 \\ 
VA & Tornado & 3 & 36.8830 & -76.0860 \\ 
\bottomrule
\end{tabular*}
\end{table}

}

\caption{\label{fig-table3}Virginia EF3 - EF5 Tornado Data}

\end{figure}%

Figure~\ref{fig-table3} shows 5 tornadoes that are only EF3 tornadoes.
This looks very similar to Florida's tornado data in
Figure~\ref{fig-table1}. The tornadoes appear to be spread out and not
concentrated in one area of the state. Below is a map of the tornadoes
for Virginia to get a better idea of where the tornadoes are located.

\begin{figure}

\centering{

\pandocbounded{\includegraphics[keepaspectratio]{results_files/figure-pdf/fig-plot4-1.pdf}}

}

\caption{\label{fig-plot4}Virginia EF3 - EF5 Tornado Map}

\end{figure}%

From Figure~\ref{fig-plot4}, there is only one tornado that could have
the potential to become a Sharknado. The tornado starts on the land so
that limits the chances of it becoming a Sharknado. This is also the
same as what happened with the Florida tornado, an EF3 tornado is really
only able to toss a car or truck (Quach (2025)). This makes Virginia an
unlikely candidate to have a Sharknado occur, diminishing the risk of an
insurance payout. EF3 tornadoes can displace the shark if its close
enough to land but that is still fairly unlikely.

\section{Results: Oklahoma}\label{results-oklahoma}

Oklahoma is known as the tornado capital of the United States. Tornado
data from Oklahoma will give a good comparison for the amount of
tornadoes that occur at the levels needed to cause a Sharknado.

\begin{figure}

\centering{

\begin{table}
\fontsize{12.0pt}{14.0pt}\selectfont
\begin{tabular*}{\linewidth}{@{\extracolsep{\fill}}llrrr}
\toprule
STATE\_ABBR & EVENT\_TYPE & TOR\_F\_SCALE & BEGIN\_LAT & BEGIN\_LON \\ 
\midrule\addlinespace[2.5pt]
OK & Tornado & 5 & 35.4440 & -98.2870 \\ 
OK & Tornado & 5 & 35.3030 & -97.6050 \\ 
OK & Tornado & 4 & 35.0080 & -97.9610 \\ 
OK & Tornado & 4 & 34.9390 & -97.6700 \\ 
OK & Tornado & 4 & 35.1890 & -97.6700 \\ 
OK & Tornado & 4 & 34.3570 & -99.1940 \\ 
OK & Tornado & 4 & 35.3080 & -97.1420 \\ 
OK & Tornado & 4 & 35.2840 & -97.6280 \\ 
OK & Tornado & 4 & 34.5590 & -97.3570 \\ 
OK & Tornado & 3 & 34.2923 & -96.2509 \\ 
\bottomrule
\end{tabular*}
\end{table}

}

\caption{\label{fig-table4}Oklahoma EF3 - EF5 Tornado Data, Ten Highest
Tornadoes}

\end{figure}%

From Figure~\ref{fig-table4}, this displays only the ten highest EF
scale tornadoes to hit Oklahoma because there are 27 tornadoes in the
EF3 to EF5 rating. The ten highest tornadoes in Oklahoma over the last
10 years have been two EF5 tornadoes, seven EF4 tornadoes, and one EF3
tornado. This data tells that EF5s are a very rare occurrence even for
the tornado capital of the country. Sharknado's would be very rare even
in a place that gets strong tornadoes at the highest rate in the
country. There is no risk of a Sharknado in Oklahoma due to being a land
locked state. There is also no insurance payout risk because of this.

\section{Results: Califorinia \& Los Angeles,
CA}\label{results-califorinia-los-angeles-ca}

To make a real connection to the Sharknado movie, examining the actual
setting of the movie, which is California and Los Angeles, CA is the
next step. The tornadoes recorded in California over the past ten years
give a real idea of the possible strength of tornadoes that could hit
Los Angeles, to see how it compares to the Sharknado film. The table
below shows the top 10 tornadoes in EF scale strength over the last 10
years.

\begin{figure}

\centering{

\begin{table}
\fontsize{12.0pt}{14.0pt}\selectfont
\begin{tabular*}{\linewidth}{@{\extracolsep{\fill}}llrrr}
\toprule
STATE\_ABBR & EVENT\_TYPE & TOR\_F\_SCALE & BEGIN\_LAT & BEGIN\_LON \\ 
\midrule\addlinespace[2.5pt]
CA & Tornado & 2 & 37.3500 & -119.3400 \\ 
CA & Tornado & 1 & 38.7050 & -121.1157 \\ 
CA & Tornado & 1 & 36.8400 & -119.5600 \\ 
CA & Tornado & 1 & 39.7313 & -120.1284 \\ 
CA & Tornado & 1 & 39.7239 & -120.1303 \\ 
CA & Tornado & 1 & 37.2500 & -119.2000 \\ 
CA & Tornado & 1 & 38.0834 & -120.7581 \\ 
CA & Tornado & 1 & 37.9300 & -120.5300 \\ 
CA & Tornado & 1 & 33.9970 & -118.1238 \\ 
CA & Tornado & 1 & 39.8043 & -121.8366 \\ 
\bottomrule
\end{tabular*}
\end{table}

}

\caption{\label{fig-table5}California Tornado Data}

\end{figure}%

Figure~\ref{fig-table5} shows the 10 highest tornadoes from the last 10
years in California. There are no tornadoes that are EF3 or higher that
have taken place in CA from this table. There are no tornadoes that can
pick up a great white shark.

Sharknado takes place in Los Angeles, CA, as the main setting for the
movie. Because of this, the table below shows only tornadoes from Los
Angeles County. Los Angeles County is where the city of Los Angeles is
located.

\begin{figure}

\centering{

\begin{table}
\fontsize{12.0pt}{14.0pt}\selectfont
\begin{tabular*}{\linewidth}{@{\extracolsep{\fill}}llrrr}
\toprule
STATE\_ABBR & EVENT\_TYPE & TOR\_F\_SCALE & BEGIN\_LAT & BEGIN\_LON \\ 
\midrule\addlinespace[2.5pt]
CA & Tornado & 0 & 33.9823 & -118.2913 \\ 
CA & Tornado & 0 & 34.5203 & -117.7595 \\ 
CA & Tornado & 0 & 33.9100 & -117.9800 \\ 
CA & Tornado & 1 & 33.9970 & -118.1238 \\ 
CA & Tornado & 0 & 33.8755 & -118.2652 \\ 
CA & Tornado & 0 & 33.8811 & -118.2153 \\ 
CA & Tornado & 0 & 33.9948 & -118.0837 \\ 
\bottomrule
\end{tabular*}
\end{table}

}

\caption{\label{fig-table6}Los Angeles, California Tornado Data}

\end{figure}%

All of the tornado data over the last ten years in Los Angeles can be
seen in Figure~\ref{fig-table6}. As shown, there have been only EF0 and
EF1 tornadoes to hit Los Angeles over this span of time. By comparing
the levels of EF ratings specified before (needing between EF3 - EF5),
we can see there is not really a possibility of getting a tornado with
wind speeds strong enough to pick up a shark to create a Sharknado. The
mph of wind speed needed to lift a shark off the ground is about 200 mph
(Salaita (n.d.)). The wind speed of an EF0 tornado is 65 - 85 mph, while
an EF1 tornado has wind speeds 86 - 110 mph. This shows us that we are
barely able to get to half the wind speed needed to pick a shark up off
the ground in Los Angeles. Below is a map showing the tornadoes that
occurred in California over the last 10 years.

\begin{figure}

\centering{

\pandocbounded{\includegraphics[keepaspectratio]{results_files/figure-pdf/fig-plot5-1.pdf}}

}

\caption{\label{fig-plot5}California Tornado Map}

\end{figure}%

As seen in Figure~\ref{fig-plot5} and Figure~\ref{fig-table5}, there are
no tornadoes strong enough to be able to lift a shark at all over the
last 10 years. There are tornadoes that have started along the coast but
each one lacked the EF strength needed to lift an object like a shark.
As a result, there is no real risk of an insurance payout for a
Sharknado in the state of California.

\section{Results: All States}\label{results-all-states}

A stacked bar chart is used to see an easier comparison of all the
tornadoes for each location we chose (Florida, South Carolina, Virginia,
Oklahoma, California) with an EF scale rating of EF3 to EF5, and compare
them to see which locations have the highest possibility of a Sharknado
to occur. Note that California only has EF0 and EF1 tornadoes, these
will be the only magnitude used in this chart.

\begin{figure}

\centering{

\pandocbounded{\includegraphics[keepaspectratio]{results_files/figure-pdf/fig-plot6-1.pdf}}

}

\caption{\label{fig-plot6}Bar Chart for Tornadoes EF3 - EF5 Over Last 10
Years for FL, SC, VA, and OK, All Tornadoes Over Last 10 Years for CA}

\end{figure}%

Figure~\ref{fig-plot6} allows us to look at the strength of each
location's tornadoes on the EF scale compared to each other. The wind
speed needed to lift a shark off the ground is 200 mph (Salaita (n.d.)).
An EF3 tornado has wind speeds of 136-165 mph, an EF4 tornado has speeds
between 166-200 mph, and an EF5 tornado has wind speeds that are over
200 mph. This means that an EF4 or EF5 tornado is only strong enough to
lift a shark off the ground. It is shown that there have only been two
tornadoes in the last ten years that have been strong enough to have
wind speeds of 200 mph or more. So there were only two tornadoes that
were strong enough to be able to pick up a shark, but they happened in
Oklahoma. That means there is no chance that there is a possibility that
a Sharknado can happen in Florida, South Carolina, or Los Angeles.

But in Oklahoma, there might be a possibility of a Sharknado. At the
Oklahoma Aquarium in Jenks, OK, there is Shark Adventure, which is home
to ten of the ``most dangerous sharks known to man'' (Aquarium (n.d.)).
These sharks are bull sharks. Jenks, OK, is in Tulsa County, so looking
at tornadoes over the last 10 years for Tulsa County, Oklahoma will give
an insight into this possibility.

\begin{figure}

\centering{

\begin{table}
\fontsize{12.0pt}{14.0pt}\selectfont
\begin{tabular*}{\linewidth}{@{\extracolsep{\fill}}llrrr}
\toprule
STATE\_ABBR & EVENT\_TYPE & TOR\_F\_SCALE & BEGIN\_LAT & BEGIN\_LON \\ 
\midrule\addlinespace[2.5pt]
OK & Tornado & 2 & 36.2180 & -96.0013 \\ 
OK & Tornado & 2 & 36.1097 & -95.9367 \\ 
OK & Tornado & 1 & 36.0926 & -96.0060 \\ 
OK & Tornado & 1 & 35.9412 & -95.8952 \\ 
OK & Tornado & 1 & 35.9792 & -95.7744 \\ 
OK & Tornado & 1 & 36.0919 & -95.7936 \\ 
OK & Tornado & 1 & 36.1742 & -95.9725 \\ 
OK & Tornado & 1 & 36.2334 & -95.9153 \\ 
OK & Tornado & 1 & 36.0058 & -96.0297 \\ 
OK & Tornado & 1 & 36.0337 & -96.0299 \\ 
\bottomrule
\end{tabular*}
\end{table}

}

\caption{\label{fig-table7}Tulsa County, OK Tornado Data}

\end{figure}%

Figure~\ref{fig-table7} shows all tornadoes that have gone through Tulsa
County, Oklahoma, and shows that these tornadoes don't reach EF3 or
higher at all. But Bull Sharks only weigh around 500 lbs ({``Bull Shark.
National Wildlife Federation''} (n.d.)). This means that they are much
easier for a tornado to lift up than a great white shark. An EF2 tornado
has the strength to slightly move houses off their foundations and
destroy mobile homes ({``The Weighty Truth: How Much Can a Tornado Lift?
- Oreate {AI} Blog''} (2025)). If an EF2 tornado went through the
Oklahoma Aquarium in Jenks and destroyed the aquarium, it could
potentially pick up sharks from the ``most dangerous sharks known to
man'' exhibit to create a real-life Sharknado. Thinking of this logic,
this could also apply to a lot of other aquariums throughout the
country. But the chances of this happening are relatively low. The last
time a tornado hit an aquarium was in 2024 when an EF1 tornado went
through the Pittsburgh Zoo \& Aquarium (Shannon (2024)). This only
resulted in little damage to the property and some trees being knocked
over. This means that an EF1 tornado going through a zoo/aquarium will
not cause enough damage to destroy buildings and have wind speeds strong
enough to pick up objects as heavy as a shark. Thinking about the
possibility of a Sharknado is pretty slim when there is a record of a
tornado that went through an aquarium; there was still no chance.

With the insurance payout possibilities in mind, storms and tornadoes
that are not shark-infested will still cause property damage. Storms and
tornadoes with high wind speeds and storms with varying precipitation
types and amounts are still possible and can still cause massive amounts
of property damage. Sharknados may be out of the picture, but there are
many other ways that weather-related events can cause property damage.

\bookmarksetup{startatroot}

\chapter{Conclusion}\label{conclusion}

In conclusion, this study shows that the likelihood of a Sharknado
occurring is very low. By comparing the weight of a great white shark to
different car types of similar weight, there is a way to determine the
strength of a tornado on the EF (Enhanced Fujita) Scale, able to pick up
a car (or a great white shark of similar weight). The EF rating required
for this to happen is an EF3-EF5 tornado. This range of tornado strength
has wind speeds of 136 mph to 200+ mph. For a Sharknado to happen,
tornadoes of this strength need to happen where great white sharks are.
Florida, South Carolina, Virginia, Califorina, and Los Angeles,
California (chosen because it's the setting of the Sharknado movie) are
coastal places. Oklahoma was also chosen to compare the amount and EF
scale types with the other places mentioned. This study looks at tornado
data from 2015-2026 and their strength from the EF scale.

Florida, South Carolina, and Virginia have very few tornadoes with the
desired strength to be able to pick up great white sharks. Florida had
just five tornadoes with an EF3 scale rating, South Carolina had nine
EF3 tornadoes and one EF4 tornado, while Virginia had just five EF3
scale rating like Florida. These tornadoes all happened inland, so the
possibility of picking up a great white shark is very unlikely.
California and Los Angeles, California have had no tornadoes of strength
EF3 or higher in the last 10 years. There will not be another Sharknado
hitting LA again.

Even though Oklahoma has had many tornadoes strong enough (EF3 to EF5)
to be able to pick up a great white shark over the last 10 years, there
are no great whites in Oklahoma. The most likely scenario to make a
Sharknado inland is having an EF3-EF5 tornado destroy an aquarium in
inland places where strong enough tornadoes occur. Even then, this event
is not very likely.

Overall, the insurance payouts for property damage from the event of a
Sharknado are extremely unlikely, if not a zero percent chance.
Tornadoes and strong storms will still exist, which is still a concern
for property damage purposes.

\bookmarksetup{startatroot}

\chapter*{References}\label{references}
\addcontentsline{toc}{chapter}{References}

\markboth{References}{References}

\phantomsection\label{refs}
\begin{CSLReferences}{1}{0}
\bibitem[\citeproctext]{ref-noauthor_shark_nodate}
Aquarium, Oklahoma. n.d. {``Shark Adventure {\textbar} Oklahoma Aquarium
Jenks, {OK}.''} Accessed March 11, 2026.
\url{https://www.okaquarium.org/176/Shark-Adventure}.

\bibitem[\citeproctext]{ref-noauthor_how_nodate}
AutoInsurance. n.d. {``How Much Does the Average Car Weigh?
{AutoInsurance}.com.''} Accessed March 9, 2026.
\url{https://www.autoinsurance.com/guide/average-car-weight/}.

\bibitem[\citeproctext]{ref-noauthor_bull_nodate}
{``Bull Shark. National Wildlife Federation.''} n.d. Accessed March 12,
2026.
\url{https://www.nwf.org/Home/Educational-Resources/Wildlife-Guide/Fish/Bull-Shark}.

\bibitem[\citeproctext]{ref-zotero-item-1}
Climatology Office, South Carolina State. 2020. {``South Carolina
Tornado Outbreak -- {OPEN} {FILE} {REPORT},''} April.

\bibitem[\citeproctext]{ref-fisheries_white_2025}
Fisheries, NOAA. 2025. {``White Shark {\textbar} {NOAA} Fisheries.
{NOAA}.''} February 19, 2025.
\url{https://www.fisheries.noaa.gov/species/white-shark}.

\bibitem[\citeproctext]{ref-quach_can_2025}
Quach, Jacqueline. 2025. {``Can a Car Be Picked up by a Tornado?
{\textbar} Otero Property Adjusting \& Appraisals.''} November 6, 2025.
\url{https://oteroadjusting.com/can-a-car-be-picked-up-by-a-tornado/}.

\bibitem[\citeproctext]{ref-salaita_how_nodate}
Salaita, Meisa. n.d. {``How a Sharknado Would Work. {HowStuffWorks}.''}
Accessed March 10, 2026.
\url{https://science.howstuffworks.com/nature/natural-disasters/sharknado.htm}.

\bibitem[\citeproctext]{ref-noauthor_pittsburgh_2024}
Shannon, Jess. 2024. {``Pittsburgh Zoo \& Aquarium Assessing Damage
After {EF}-1 Tornado Touched down in Highland Park. Yahoo News.''} May
17, 2024.
\url{https://www.yahoo.com/news/pittsburgh-zoo-aquarium-assessing-damage-233609618.html}.

\bibitem[\citeproctext]{ref-SharknadoTVMovie}
{``Sharknado ({TV} Movie 2013) - Filming \& Production - {IMDb}.''}
2013. July 11, 2013.
\url{https://www.imdb.com/title/tt2724064/locations/}.

\bibitem[\citeproctext]{ref-StormEventsDatabaseLA}
{``Storm Events Database - Search Results California National Centers
for Environmental Information.''} n.d. Accessed February 20, 2026.
\url{https://www.ncei.noaa.gov/stormevents/listevents.jsp?eventType=\%28C\%29+Tornado&beginDate_mm=10&beginDate_dd=01&beginDate_yyyy=2010&endDate_mm=10&endDate_dd=31&endDate_yyyy=2025&county=LOS\%2BANGELES\%3A37&hailfilter=0.00&tornfilter=0&windfilter=000&sort=DT&submitbutton=Search&statefips=6\%2CCALIFORNIA}.

\bibitem[\citeproctext]{ref-StormEventsDatabaseFL}
{``Storm Events Database - Search Results Florida National Centers for
Environmental Information.''} n.d. Accessed February 20, 2026.
\url{https://www.ncei.noaa.gov/stormevents/listevents.jsp?eventType=\%28C\%29+Tornado&beginDate_mm=10&beginDate_dd=01&beginDate_yyyy=2015&endDate_mm=10&endDate_dd=31&endDate_yyyy=2025&county=ALL&hailfilter=0.00&tornfilter=0&windfilter=000&sort=DT&submitbutton=Search&statefips=12\%2CFLORIDA}.

\bibitem[\citeproctext]{ref-StormEventsDatabaseCA}
{``Storm Events Database - Search Results {\textbar} National Centers
for Environmental Information.''} n.d.a. Accessed April 8, 2026.
\url{https://www.ncei.noaa.gov/stormevents/listevents.jsp?eventType=\%28C\%29+Tornado&beginDate_mm=12&beginDate_dd=01&beginDate_yyyy=2015&endDate_mm=12&endDate_dd=31&endDate_yyyy=2025&county=ALL&hailfilter=0.00&tornfilter=0&windfilter=000&sort=DT&submitbutton=Search&statefips=6\%2CCALIFORNIA}.

\bibitem[\citeproctext]{ref-StormEventsDatabaseVA}
---------. n.d.b. Accessed April 14, 2026.
\url{https://www.ncei.noaa.gov/stormevents/listevents.jsp?eventType=\%28C\%29+Tornado&beginDate_mm=12&beginDate_dd=01&beginDate_yyyy=2015&endDate_mm=12&endDate_dd=31&endDate_yyyy=2025&county=ALL&hailfilter=0.00&tornfilter=3&windfilter=000&sort=DT&submitbutton=Search&statefips=51\%2CVIRGINIA}.

\bibitem[\citeproctext]{ref-StormEventsDatabaseOK}
{``Storm Events Database - Search Results Oklahoma National Centers for
Environmental Information.''} n.d. Accessed February 20, 2026.
\url{https://www.ncei.noaa.gov/stormevents/listevents.jsp?eventType=\%28C\%29+Tornado&beginDate_mm=10&beginDate_dd=01&beginDate_yyyy=2010&endDate_mm=10&endDate_dd=31&endDate_yyyy=2025&county=ALL&hailfilter=0.00&tornfilter=0&windfilter=000&sort=DT&submitbutton=Search&statefips=40\%2COKLAHOMA}.

\bibitem[\citeproctext]{ref-StormEventsDatabaseSC}
{``Storm Events Database - Search Results South Carolina National
Centers for Environmental Information.''} n.d. Accessed February 20,
2026.
\url{https://www.ncei.noaa.gov/stormevents/listevents.jsp?eventType=\%28C\%29+Tornado&beginDate_mm=10&beginDate_dd=01&beginDate_yyyy=2010&endDate_mm=10&endDate_dd=31&endDate_yyyy=2025&county=ALL&hailfilter=0.00&tornfilter=0&windfilter=000&sort=DT&submitbutton=Search&statefips=45\%2CSOUTH+CAROLINA\#}.

\bibitem[\citeproctext]{ref-NOAA}
{``Storm Events Database {\textbar} National Centers for Environmental
Information.''} n.d. Accessed April 22, 2026.
\url{https://www.ncei.noaa.gov/stormevents/}.

\bibitem[\citeproctext]{ref-noauthor_weighty_2025}
{``The Weighty Truth: How Much Can a Tornado Lift? - Oreate {AI}
Blog.''} 2025. December 29, 2025.
\url{https://www.oreateai.com/blog/the-weighty-truth-how-much-can-a-tornado-lift/22e09faa7969386318853dd21fb2fbe4}.

\bibitem[\citeproctext]{ref-us_department_of_commerce_enhanced_nodate}
US Department of Commerce, NOAA. n.d. {``The Enhanced Fujita Scale ({EF}
Scale).''} {NOAA}'s National Weather Service. Accessed March 11, 2026.
\url{https://www.weather.gov/oun/efscale}.

\end{CSLReferences}

\cleardoublepage
\phantomsection
\addcontentsline{toc}{part}{Appendices}
\appendix

\chapter{Draft: Data Documentation}\label{draft-data-documentation}

\begin{tcolorbox}[enhanced jigsaw, coltitle=black, leftrule=.75mm, colframe=quarto-callout-note-color-frame, bottomtitle=1mm, bottomrule=.15mm, arc=.35mm, colback=white, colbacktitle=quarto-callout-note-color!10!white, titlerule=0mm, opacitybacktitle=0.6, breakable, toptitle=1mm, title=\textcolor{quarto-callout-note-color}{\faInfo}\hspace{0.5em}{Note}, left=2mm, toprule=.15mm, rightrule=.15mm, opacityback=0]

Put your name on your work and give it an appropriate title -- edit your
\_quarto.yml file

This needs to be in paragraph format (notes are fine in the draft, but
for the real thing, please format everything and connect things with
words.)

Make sure you add a date -- the release date? -- to your bibliography
entry, and make sure the other fields are filled in. Zotero doesn't
always catch everything perfectly.

\end{tcolorbox}

\chapter{Data Documentation}\label{data-documentation-1}

All data was collected from NOAA (National Oceanic and Atmospheric
Administration) Storm Events Database. The NOAA's storm events database
contains records on various types fo severe weather collected by NOAA's
National Weather Service.

Tornado data from NOAA Storm Events Database was collected for five
different states and one county across the United States ({``Storm
Events Database {\textbar} National Centers for Environmental
Information''} (n.d.)). These locations were chosen because of their
amount of EF3 or higher tornadoes, location close to the ocean, and
setting of the Sharknado movie. These locations are:

\begin{itemize}
\item
  Florida: high frequency of tornadoes EF3 or higher for a coastal
  region
\item
  South Carolina: high frequency of tornadoes EF3 or higher for a
  coastal region
\item
  Virginia: high frequency of tornadoes EF3 or higher for a coastal
  region
\item
  Oklahoma: high amount of powerful tornadoes to compare the frequency
  of high magnitude tornadoes
\item
  California and Los Angeles, CA: State and city setting of the movie
  ``Sharknado'' ({``Sharknado ({TV} Movie 2013) - Filming \& Production
  - {IMDb}''} (2013))
\end{itemize}

All data was analyzed using the programming software R and these
following R packages:

\begin{itemize}
\item
  \texttt{tidyverse}: data cleaning, data manipulation, and table
  formulation
\item
  \texttt{gt}: to make better looking tables
\item
  \texttt{maps}: make maps of United States for tornado data
\end{itemize}

\begin{tcolorbox}[enhanced jigsaw, coltitle=black, leftrule=.75mm, colframe=quarto-callout-note-color-frame, bottomtitle=1mm, bottomrule=.15mm, arc=.35mm, colback=white, colbacktitle=quarto-callout-note-color!10!white, titlerule=0mm, opacitybacktitle=0.6, breakable, toptitle=1mm, title=\textcolor{quarto-callout-note-color}{\faInfo}\hspace{0.5em}{Note}, left=2mm, toprule=.15mm, rightrule=.15mm, opacityback=0]

These are fine (other than the author being missing), but you need to
incorporate them into your writing as an organic part of the paragraphs.

\end{tcolorbox}

\chapter{Draft: Results}\label{draft-results}

\chapter{Methods}\label{methods-2}

A Sharknado event involves a tornado being able to lift a great white
shark out of the ocean and carry it onto land. To see if this is
possible, we need to look at the physics: the weight of a great white
shark, what objects are known to be picked up by tornadoes, and the
magnitude of tornadoes required to pick up a shark. First, finding out
how a great white shark could be picked up by a tornado is nessesary.
The weight of a great white shark will be compared to other objects that
can and have been picked up by tornadoes, such as cars. Next, figuring
out what strength (in EF scale from EF3 - EF5) of a tornado would be
needed to be able to lift a great white shark and a car. Lastly,
identifying tornadoes from 2015-2026 that are in that EF range from 6
different places in the United States. These places are be Florida,
South Carolina, Virginia, Oklahoma, California, and Los Angeles. These
places were chosen for specific reasons. Florida has many tropical
storms and is a coastal state (Florida, (n.d.)). South Carolina is also
another place that can be hit by tropical storms and is another coastal
state (South Carolina, (n.d.)). Virginia is a coastal state that has had
a good amount of tornadoes of our criteria (Virginia, (n.d.b)). Oklahoma
is not a coastal state, but it is the tornado capital of the United
States, so we will just look at the number of tornadoes that meet the
criteria to get an idea of how rare strong tornadoes are (Oklahoma,
(n.d.)). California ((n.d.a)) was chosen because it is the state setting
of Sharknado, while Los Angeles ((n.d.)) was picked because it is the
city setting of the Sharknado movie ({``Sharknado ({TV} Movie 2013) -
Filming \& Production - {IMDb}''} (2013)).

Sharks and cars weigh relatively the same in pounds (lbs). The average
great white shark weighs around 4,500 lbs (Fisheries (2025)). A car SUV
weighs around 3,986 lbs while a truck SUV weighs around 4,578 lbs
(AutoInsurance (n.d.)). This means that your average great white shark
weighs more than a big car and as much as a truck or SUV. We will use a
car as the base to find what EF scale level will be used to lift a great
white shark out of the water.

The EF scale is made up of 6 different ratings of tornado based on
3-second wind speed in mph (US Department of Commerce (n.d.)). An EF0
tornado has the lowest wind speed, ranging between 65-85 mph, EF1 has
wind speeds between 86-110 mph, EF2 has wind speeds between 111-135 mph,
EF3 has wind speeds of 136-165 mph, EF4 has speeds between 166-200 mph,
and EF5 has wind speeds that are over 200 mph. An EF3 tornado is likely
to toss cars and trucks, while an EF4 is very likely to toss cars many
feet, and an EF5 tornado is almost certain to lift and destroy cars. EF3
tornadoes are unlikely to fully pick up a shark but it is still
important to analyze them because a great white shark could be moved by
an EF3 tornado. Because of the comparison of a great white shark to the
size and weight of a car SUV, and a truck SUV, that gives a good
comparison to be able to put a shark in the place of such cars in the EF
scales' ability to lift cars (Quach (2025)). Tornadoes from EF3 to EF5
are chosen because this is the strength of tornadoes that are able to
lift a car (Salaita (n.d.)).

Using these metrics, finding storms that have had the potential to
produce a Sharknado allows the ability to evaluate the potential
insurance risk with such events.

\chapter{Results}\label{results-1}

The results of this study examines tornado activity over the last 10
years from specific states across the United States. Tornado strength of
EF3 - EF5 is chosen due to their capabilities to lift and toss objects
that are similar to sharks weight wise, these objects being cars or
trucks. The findings are shown by state and area (FL, SC, VA, CA, Los
Angeles) to be able to see where and which tornadoes had the strength
and location to start a Sharknado event. California and Los Angeles, CA
were chosen due to being the location of the film ``Sharknado''.

Below are these tornadoes being mapped on the entire United States to
get an idea of where these states and their tornadoes are and why they
were chosen for this study.

\begin{figure}

\centering{

\pandocbounded{\includegraphics[keepaspectratio]{draft-results_files/figure-pdf/fig-plot1-1.pdf}}

}

\caption{\label{fig-plot1}EF3 - EF5 Tornado Data Accross USA}

\end{figure}%

Most tornadoes happen on the eastern side of the country as shown in
Figure~\ref{fig-plot1}. This shows the regions to check to see if
Sharknadoes are possible. The east coast is the best place to start the
search. The states that have enough tornadoes, and the states that will
be examined, are Florida, South Carolina, and Virginia.

\section{Results: Florida}\label{results-florida-1}

From Figure~\ref{fig-plot1}, Florida is a place that is on the coast.
This would make it possible to pick up a shark in the event of a strong
enough tornado. There are enough tornadoes that touched down in Florida
over the last 10 years which is a reason Florida is a state that will be
analyzed.

\begin{figure}

\centering{

\begin{table}
\fontsize{12.0pt}{14.0pt}\selectfont
\begin{tabular*}{\linewidth}{@{\extracolsep{\fill}}llrrr}
\toprule
STATE\_ABBR & EVENT\_TYPE & TOR\_F\_SCALE & BEGIN\_LAT & BEGIN\_LON \\ 
\midrule\addlinespace[2.5pt]
FL & Tornado & 3 & 30.8707 & -87.3810 \\ 
FL & Tornado & 3 & 30.4907 & -87.2052 \\ 
FL & Tornado & 3 & 30.4934 & -84.1262 \\ 
FL & Tornado & 3 & 30.6376 & -85.5108 \\ 
FL & Tornado & 3 & 30.1384 & -85.7526 \\ 
FL & Tornado & 3 & 26.8878 & -81.0381 \\ 
FL & Tornado & 3 & 26.5555 & -80.3321 \\ 
FL & Tornado & 3 & 27.3740 & -80.4270 \\ 
\bottomrule
\end{tabular*}
\end{table}

}

\caption{\label{fig-table1}Florida EF3 - EF5 Tornado Data}

\end{figure}%

Figure~\ref{fig-table1} shows that there have been only 5 EF3 tornadoes
in the last 10 years in the state of Florida. Because of an EF3 tornado
being the highest EF rating, and it being ``likely to toss a car or
truck'', a tornado in Florida is fairly unlikely to pick up the average
great white shark. Having only five EF3 tornadoes in the past 10 years
seems like the chances of a Sharknado happening seem pretty slim. Below
is a map showing these tornadoes.

\begin{figure}

\centering{

\pandocbounded{\includegraphics[keepaspectratio]{draft-results_files/figure-pdf/fig-plot2-1.pdf}}

}

\caption{\label{fig-plot2}Florida EF3 - EF5 Tornado Map}

\end{figure}%

Figure~\ref{fig-plot2} shows a map of Florida showing the five EF3
tornadoes that have occurred in the last 10 years. There is one tornado
that shows up close to the water, but looking at the coordinates, it
appears on the beach. EF3 tornadoes are likely to toss cars and trucks
(Quach (2025)). Tossing and picking up two different things in this
scenario, leading to believe an EF3 tornado is not very likely to fully
pick up a great white shark. The EF3 tornadoes can displace the shark if
its close enough to land but that is still fairly unlikely.

\section{Results: South Carolina}\label{results-south-carolina-1}

A place with similar properties to Florida, having tropical storms
frequently and being a coastal state, is South Carolina. The data below
gives another state to compare to that is in an area where great white
sharks may be present.

\begin{figure}

\centering{

\begin{table}
\fontsize{12.0pt}{14.0pt}\selectfont
\begin{tabular*}{\linewidth}{@{\extracolsep{\fill}}llrrr}
\toprule
STATE\_ABBR & EVENT\_TYPE & TOR\_F\_SCALE & BEGIN\_LAT & BEGIN\_LON \\ 
\midrule\addlinespace[2.5pt]
SC & Tornado & 3 & 34.6170 & -83.0840 \\ 
SC & Tornado & 3 & 33.2751 & -81.7048 \\ 
SC & Tornado & 3 & 33.2754 & -81.5438 \\ 
SC & Tornado & 3 & 33.3177 & -81.2903 \\ 
SC & Tornado & 3 & 33.2071 & -81.2742 \\ 
SC & Tornado & 3 & 33.4539 & -81.2651 \\ 
SC & Tornado & 4 & 32.7045 & -81.2899 \\ 
SC & Tornado & 3 & 33.1816 & -79.9901 \\ 
SC & Tornado & 3 & 32.9286 & -81.3756 \\ 
SC & Tornado & 3 & 33.1154 & -81.1889 \\ 
\bottomrule
\end{tabular*}
\end{table}

}

\caption{\label{fig-table2}South Carolina EF3 - EF5 Tornado Data}

\end{figure}%

South Carolina is similar to Florida from Figure~\ref{fig-table2}. There
are a lot of EF3 tornadoes with one EF4 tornado. But it's in the same
boat as Florida; there has been a very limited number of powerful
tornadoes over the past 10 years in South Carolina. There have been nine
EF3 tornadoes and one EF4 tornado over that span. The surprising thing
is that the EF4 and six of the EF3 tornadoes all happened on the same
day. But the possibility of a Sharknado is very low because these
tornadoes all happened away from the coast. The EF4 happened in Estill
and Hampton Counties (Climatology Office (2020)). Below is a map showing
all the tornadoes from South Carolina over the last 10 years.

\begin{figure}

\centering{

\pandocbounded{\includegraphics[keepaspectratio]{draft-results_files/figure-pdf/fig-plot3-1.pdf}}

}

\caption{\label{fig-plot3}South Carolina EF3 - EF5 Tornado Map}

\end{figure}%

The EF4 tornado is the main Sharknado threat. As seen in
Figure~\ref{fig-plot3}, all tornadoes EF3 and above only appear on land
and are far enough from the coast to have any real threat of a
Sharknado. From the evidence from Florida and now South Carolina, a
Sharknado is looking less and less possible. The insurance payout risk
associated with a Sharknado in South Carolina is negligible.

\section{Results: Virginia}\label{results-virginia-1}

Another coastal state similar to South Carolina is Virginia.
Figure~\ref{fig-plot1} shows a plethora of tornadoes in VA over the last
10 years that fit the criteria of EF3 or higher tornadoes to test the
potential of a Sharknado. Below is Virginia tornado data from the last
10 years.

\begin{figure}

\centering{

\begin{table}
\fontsize{12.0pt}{14.0pt}\selectfont
\begin{tabular*}{\linewidth}{@{\extracolsep{\fill}}llrrr}
\toprule
STATE\_ABBR & EVENT\_TYPE & TOR\_F\_SCALE & BEGIN\_LAT & BEGIN\_LON \\ 
\midrule\addlinespace[2.5pt]
VA & Tornado & 3 & 37.2318 & -78.8576 \\ 
VA & Tornado & 3 & 37.8285 & -76.9681 \\ 
VA & Tornado & 3 & 37.4586 & -79.1838 \\ 
VA & Tornado & 3 & 36.8564 & -79.8906 \\ 
VA & Tornado & 3 & 36.8830 & -76.0860 \\ 
\bottomrule
\end{tabular*}
\end{table}

}

\caption{\label{fig-table3}Virginia EF3 - EF5 Tornado Data}

\end{figure}%

Figure~\ref{fig-table3} shows 5 tornadoes that are only EF3 tornadoes.
This looks very similar to Florida's tornado data in
Figure~\ref{fig-table1}. The tornadoes appear to be spread out and not
concentrated in one area of the state. Below is a map of the tornadoes
for Virginia to get a better idea of where the tornadoes are located.

\begin{figure}

\centering{

\pandocbounded{\includegraphics[keepaspectratio]{draft-results_files/figure-pdf/fig-plot4-1.pdf}}

}

\caption{\label{fig-plot4}Virginia EF3 - EF5 Tornado Map}

\end{figure}%

From Figure~\ref{fig-plot4}, there is only one tornado that could have
the potential to become a Sharknado. The tornado starts on the land so
that limits the chances of it becoming a Sharknado. This is also the
same as what happened with the Florida tornado, an EF3 tornado is really
only able to toss a car or truck (Quach (2025)). This makes Virginia an
unlikely candidate to have a Sharknado occur, diminishing the risk of an
insurance payout. EF3 tornadoes can displace the shark if its close
enough to land but that is still fairly unlikely.

\section{Results: Oklahoma}\label{results-oklahoma-1}

Oklahoma is known as the tornado capital of the United States. Tornado
data from Oklahoma will give a good comparison for the amount of
tornadoes that occur at the levels needed to cause a Sharknado.

\begin{figure}

\centering{

\begin{table}
\fontsize{12.0pt}{14.0pt}\selectfont
\begin{tabular*}{\linewidth}{@{\extracolsep{\fill}}llrrr}
\toprule
STATE\_ABBR & EVENT\_TYPE & TOR\_F\_SCALE & BEGIN\_LAT & BEGIN\_LON \\ 
\midrule\addlinespace[2.5pt]
OK & Tornado & 5 & 35.4440 & -98.2870 \\ 
OK & Tornado & 5 & 35.3030 & -97.6050 \\ 
OK & Tornado & 4 & 35.0080 & -97.9610 \\ 
OK & Tornado & 4 & 34.9390 & -97.6700 \\ 
OK & Tornado & 4 & 35.1890 & -97.6700 \\ 
OK & Tornado & 4 & 34.3570 & -99.1940 \\ 
OK & Tornado & 4 & 35.3080 & -97.1420 \\ 
OK & Tornado & 4 & 35.2840 & -97.6280 \\ 
OK & Tornado & 4 & 34.5590 & -97.3570 \\ 
OK & Tornado & 3 & 34.2923 & -96.2509 \\ 
\bottomrule
\end{tabular*}
\end{table}

}

\caption{\label{fig-table4}Oklahoma EF3 - EF5 Tornado Data, Ten Highest
Tornadoes}

\end{figure}%

From Figure~\ref{fig-table4}, this displays only the ten highest EF
scale tornadoes to hit Oklahoma because there are 27 tornadoes in the
EF3 to EF5 rating. The ten highest tornadoes in Oklahoma over the last
10 years have been two EF5 tornadoes, seven EF4 tornadoes, and one EF3
tornado. This data tells that EF5s are a very rare occurrence even for
the tornado capital of the country. Sharknado's would be very rare even
in a place that gets strong tornadoes at the highest rate in the
country. There is no risk of a Sharknado in Oklahoma due to being a land
locked state. There is also no insurance payout risk because of this.

\section{Results: Califorinia \& Los Angeles,
CA}\label{results-califorinia-los-angeles-ca-1}

To make a real connection to the Sharknado movie, examining the actual
setting of the movie, which is California and Los Angeles, CA is the
next step. The tornadoes recorded in California over the past ten years
give a real idea of the possible strength of tornadoes that could hit
Los Angeles, to see how it compares to the Sharknado film. The table
below shows the top 10 tornadoes in EF scale strength over the last 10
years.

\begin{figure}

\centering{

\begin{table}
\fontsize{12.0pt}{14.0pt}\selectfont
\begin{tabular*}{\linewidth}{@{\extracolsep{\fill}}llrrr}
\toprule
STATE\_ABBR & EVENT\_TYPE & TOR\_F\_SCALE & BEGIN\_LAT & BEGIN\_LON \\ 
\midrule\addlinespace[2.5pt]
CA & Tornado & 2 & 37.3500 & -119.3400 \\ 
CA & Tornado & 1 & 38.7050 & -121.1157 \\ 
CA & Tornado & 1 & 36.8400 & -119.5600 \\ 
CA & Tornado & 1 & 39.7313 & -120.1284 \\ 
CA & Tornado & 1 & 39.7239 & -120.1303 \\ 
CA & Tornado & 1 & 37.2500 & -119.2000 \\ 
CA & Tornado & 1 & 38.0834 & -120.7581 \\ 
CA & Tornado & 1 & 37.9300 & -120.5300 \\ 
CA & Tornado & 1 & 33.9970 & -118.1238 \\ 
CA & Tornado & 1 & 39.8043 & -121.8366 \\ 
\bottomrule
\end{tabular*}
\end{table}

}

\caption{\label{fig-table5}California Tornado Data}

\end{figure}%

Figure~\ref{fig-table5} shows the 10 highest tornadoes from the last 10
years in California. There are no tornadoes that are EF3 or higher that
have taken place in CA from this table. There are no tornadoes that can
pick up a great white shark.

Sharknado takes place in Los Angeles, CA, as the main setting for the
movie. Because of this, the table below shows only tornadoes from Los
Angeles County. Los Angeles County is where the city of Los Angeles is
located.

\begin{figure}

\centering{

\begin{table}
\fontsize{12.0pt}{14.0pt}\selectfont
\begin{tabular*}{\linewidth}{@{\extracolsep{\fill}}llrrr}
\toprule
STATE\_ABBR & EVENT\_TYPE & TOR\_F\_SCALE & BEGIN\_LAT & BEGIN\_LON \\ 
\midrule\addlinespace[2.5pt]
CA & Tornado & 0 & 33.9823 & -118.2913 \\ 
CA & Tornado & 0 & 34.5203 & -117.7595 \\ 
CA & Tornado & 0 & 33.9100 & -117.9800 \\ 
CA & Tornado & 1 & 33.9970 & -118.1238 \\ 
CA & Tornado & 0 & 33.8755 & -118.2652 \\ 
CA & Tornado & 0 & 33.8811 & -118.2153 \\ 
CA & Tornado & 0 & 33.9948 & -118.0837 \\ 
\bottomrule
\end{tabular*}
\end{table}

}

\caption{\label{fig-table6}Los Angeles, California Tornado Data}

\end{figure}%

All of the tornado data over the last ten years in Los Angeles can be
seen in Figure~\ref{fig-table6}. As shown, there have been only EF0 and
EF1 tornadoes to hit Los Angeles over this span of time. By comparing
the levels of EF ratings specified before (needing between EF3 - EF5),
we can see there is not really a possibility of getting a tornado with
wind speeds strong enough to pick up a shark to create a Sharknado. The
mph of wind speed needed to lift a shark off the ground is about 200 mph
(Salaita (n.d.)). The wind speed of an EF0 tornado is 65 - 85 mph, while
an EF1 tornado has wind speeds 86 - 110 mph. This shows us that we are
barely able to get to half the wind speed needed to pick a shark up off
the ground in Los Angeles. Below is a map showing the tornadoes that
occurred in California over the last 10 years.

\begin{figure}

\centering{

\pandocbounded{\includegraphics[keepaspectratio]{draft-results_files/figure-pdf/fig-plot5-1.pdf}}

}

\caption{\label{fig-plot5}California Tornado Map}

\end{figure}%

As seen in Figure~\ref{fig-plot5} and Figure~\ref{fig-table5}, there are
no tornadoes strong enough to be able to lift a shark at all over the
last 10 years. There are tornadoes that have started along the coast but
each one lacked the EF strength needed to lift an object like a shark.
As a result, there is no real risk of an insurance payout for a
Sharknado in the state of California.

\section{Results: All States}\label{results-all-states-1}

A stacked bar chart is used to see an easier comparison of all the
tornadoes for each location we chose (Florida, South Carolina, Virginia,
Oklahoma, California) with an EF scale rating of EF3 to EF5, and compare
them to see which locations have the highest possibility of a Sharknado
to occur. Note that California only has EF0 and EF1 tornadoes, these
will be the only magnitude used in this chart.

\begin{figure}

\centering{

\pandocbounded{\includegraphics[keepaspectratio]{draft-results_files/figure-pdf/fig-plot6-1.pdf}}

}

\caption{\label{fig-plot6}Bar Chart for Tornadoes EF3 - EF5 Over Last 10
Years for FL, SC, VA, and OK, All Tornadoes Over Last 10 Years for CA}

\end{figure}%

Figure~\ref{fig-plot6} allows us to look at the strength of each
location's tornadoes on the EF scale compared to each other. The wind
speed needed to lift a shark off the ground is 200 mph (Salaita (n.d.)).
An EF3 tornado has wind speeds of 136-165 mph, an EF4 tornado has speeds
between 166-200 mph, and an EF5 tornado has wind speeds that are over
200 mph. This means that an EF4 or EF5 tornado is only strong enough to
lift a shark off the ground. It is shown that there have only been two
tornadoes in the last ten years that have been strong enough to have
wind speeds of 200 mph or more. So there were only two tornadoes that
were strong enough to be able to pick up a shark, but they happened in
Oklahoma. That means there is no chance that there is a possibility that
a Sharknado can happen in Florida, South Carolina, or Los Angeles.

But in Oklahoma, there might be a possibility of a Sharknado. At the
Oklahoma Aquarium in Jenks, OK, there is Shark Adventure, which is home
to ten of the ``most dangerous sharks known to man'' (Aquarium (n.d.)).
These sharks are bull sharks. Jenks, OK, is in Tulsa County, so looking
at tornadoes over the last 10 years for Tulsa County, Oklahoma will give
an insight into this possibility.

\begin{figure}

\centering{

\begin{table}
\fontsize{12.0pt}{14.0pt}\selectfont
\begin{tabular*}{\linewidth}{@{\extracolsep{\fill}}llrrr}
\toprule
STATE\_ABBR & EVENT\_TYPE & TOR\_F\_SCALE & BEGIN\_LAT & BEGIN\_LON \\ 
\midrule\addlinespace[2.5pt]
OK & Tornado & 2 & 36.2180 & -96.0013 \\ 
OK & Tornado & 2 & 36.1097 & -95.9367 \\ 
OK & Tornado & 1 & 36.0926 & -96.0060 \\ 
OK & Tornado & 1 & 35.9412 & -95.8952 \\ 
OK & Tornado & 1 & 35.9792 & -95.7744 \\ 
OK & Tornado & 1 & 36.0919 & -95.7936 \\ 
OK & Tornado & 1 & 36.1742 & -95.9725 \\ 
OK & Tornado & 1 & 36.2334 & -95.9153 \\ 
OK & Tornado & 1 & 36.0058 & -96.0297 \\ 
OK & Tornado & 1 & 36.0337 & -96.0299 \\ 
\bottomrule
\end{tabular*}
\end{table}

}

\caption{\label{fig-table7}Tulsa County, OK Tornado Data}

\end{figure}%

Figure~\ref{fig-table7} shows all tornadoes that have gone through Tulsa
County, Oklahoma, and shows that these tornadoes don't reach EF3 or
higher at all. But Bull Sharks only weigh around 500 lbs ({``Bull Shark.
National Wildlife Federation''} (n.d.)). This means that they are much
easier for a tornado to lift up than a great white shark. An EF2 tornado
has the strength to slightly move houses off their foundations and
destroy mobile homes ({``The Weighty Truth: How Much Can a Tornado Lift?
- Oreate {AI} Blog''} (2025)). If an EF2 tornado went through the
Oklahoma Aquarium in Jenks and destroyed the aquarium, it could
potentially pick up sharks from the ``most dangerous sharks known to
man'' exhibit to create a real-life Sharknado. Thinking of this logic,
this could also apply to a lot of other aquariums throughout the
country. But the chances of this happening are relatively low. The last
time a tornado hit an aquarium was in 2024 when an EF1 tornado went
through the Pittsburgh Zoo \& Aquarium (Shannon (2024)). This only
resulted in little damage to the property and some trees being knocked
over. This means that an EF1 tornado going through a zoo/aquarium will
not cause enough damage to destroy buildings and have wind speeds strong
enough to pick up objects as heavy as a shark. Thinking about the
possibility of a Sharknado is pretty slim when there is a record of a
tornado that went through an aquarium; there was still no chance.

With the insurance payout possibilities in mind, storms and tornadoes
that are not shark-infested will still cause property damage. Storms and
tornadoes with high wind speeds and storms with varying precipitation
types and amounts are still possible and can still cause massive amounts
of property damage. Sharknados may be out of the picture, but there are
many other ways that weather-related events can cause property damage.

\chapter{Draft: Intro/Conclusions}\label{draft-introconclusions}

\chapter{1. Introduction}\label{introduction-1}

This reports evaluates the potential insurance risk associated with a
Sharknado scenario for property damage payouts . A Sharknado is
explained by a tornado that generates enough wind speed to be able to
lift a heavy animal, such as a great white shark, out of the water and
into its vortex. A Sharknado can also be called a ``Shark-infested
tornado''.

This study starts by comparing the average great white shark to objects
of similar weight, specifically cars or pickup trucks. Tornadoes are
known for destroying homes and tossing around vehicles, thus finding the
types of vehicles that are similar enough in weight to a great white
shark makes finding the right magnitude of tornadoes required to lift a
great white shark easier.

Next, finding the magnitude of a tornado capable of picking up a great
white shark, or a car, is found using the EF scale (Enhanced Fujita
Scale) by looking at EF ratings (US Department of Commerce (n.d.)). The
Enhanced Fujita Scale uses ratings of EF0 to EF5. Tornadoes of interest
for this study are tornadoes of rating EF3 to EF5. These tornadoes have
strong enough wind speeds to potentially pick up and throw cars or, in
this case, a great white shark.

Finally, looking at places that are near the coast to be able to pick up
a great white shark is essential to actually generating a Sharknado.
There are 4 locations observed for their tornado data that fit this
agenda in this study: Florida, South Carolina, Virginia, California, and
Los Angeles (this is the setting of the Sharknado movie). For another
comparison of tornado strength and amount, Oklahoma is also looked at
for its tornado EF rating, even though it is landlocked and does not
have any native shark population. These locations have tornado data
taken from 2015 - 2026 using NOAA tornado data.

Combining the magnitude of tornadoes with the areas where these
tornadoes touch down gives an idea of the possibility of a potential
Sharknado event occurring.

\chapter{4. Conclusion}\label{conclusion-1}

In conclusion, this study shows that the likelihood of a Sharknado
occurring is very low. By comparing the weight of a great white shark to
different car types of similar weight, there is a way to determine the
strength of a tornado on the EF (Enhanced Fujita) Scale, able to pick up
a car (or a great white shark of similar weight). The EF rating required
for this to happen is an EF3-EF5 tornado. This range of tornado strength
has wind speeds of 136 mph to 200+ mph. For a Sharknado to happen,
tornadoes of this strength need to happen where great white sharks are.
Florida, South Carolina, Virginia, Califorina, and Los Angeles,
California (chosen because it's the setting of the Sharknado movie) are
coastal places. Oklahoma was also chosen to compare the amount and EF
scale types with the other places mentioned. This study looks at tornado
data from 2015-2026 and their strength from the EF scale.

Florida, South Carolina, and Virginia have very few tornadoes with the
desired strength to be able to pick up great white sharks. Florida had
just five tornadoes with an EF3 scale rating, South Carolina had nine
EF3 tornadoes and one EF4 tornado, while Virginia had just five EF3
scale rating like Florida. These tornadoes all happened inland, so the
possibility of picking up a great white shark is very unlikely.
California and Los Angeles, California have had no tornadoes of strength
EF3 or higher in the last 10 years. There will not be another Sharknado
hitting LA again.

Even though Oklahoma has had many tornadoes strong enough (EF3 to EF5)
to be able to pick up a great white shark over the last 10 years, there
are no great whites in Oklahoma. The most likely scenario to make a
Sharknado inland is having an EF3-EF5 tornado destroy an aquarium in
inland places where strong enough tornadoes occur. Even then, this event
is not very likely.

Overall, the insurance payouts for property damage from the event of a
Sharknado are extremely unlikely, if not a zero percent chance.
Tornadoes and strong storms will still exist, which is still a concern
for property damage purposes.

\chapter{Draft: Executive Summary}\label{draft-executive-summary}

\chapter{Executive Summary}\label{executive-summary-1}

\subsection{Purpose of Report}\label{purpose-of-report-1}

This study dives into the possibility of a Sharknado event. A Sharknado
can be explained as a ``shark-infested tornado''. This is a tornado that
generates enough wind speed to lift a great white shark from the water
and into its vortex. The purpose of this study is to assess whether a
Sharknado event could be considered a risk for property damage insurance
payouts.

\subsection{Methods}\label{methods-3}

To start, a comparison of the weight of a great white shark (around 4500
lbs) and vehicles (around 4000 - 4600 lbs for SUVs and Trucks) is used
because cars are known to be lifted and destroyed by tornadoes. Using
the EF (Enhanced Fujita) Scale, the strength of tornadoes able to lift
and toss cars (or a great white shark) is found to be between EF3 - EF5
tornadoes. These are wind speeds of 136 mph to 200+ mph. This study will
only be looking at tornadoes of EF3 - EF5 for this reason. This study
also looks at four locations of tornado data from 2015 to 2026. These
locations are Florida, South Carolina, Virginia, Oklahoma, California,
and Los Angeles County, CA. These places were chosen because of their
abundant tropical storms and coastal area (Florida and South Carolina),
historic tornado amounts and strength (Oklahoma), and the setting of the
movie Sharknado (Los Angeles, CA).

\subsection{Findings}\label{findings-1}

Tornadoes of strength EF3 - EF5 are fairly uncommon for coastal states
like Florida and South Carolina. Florida had only five tornadoes in the
last 10 years that were of strength EF3, and zero EF4 and EF5 tornadoes.
All but one of these tornadoes started and ended on land. The one
tornado started on a beach in Florida, where the water is not deep
enough for sharks. South Carolina had nine EF3 and one EF4 tornado over
the last 10 years. All of these tornadoes started and stayed on land,
making a Sharknado impossible. Virginia had five EF3 tornadoes that all
happened on land. Oklahoma has had plenty of tornadoes strong enough to
lift and carry a shark, but there are no sharks in Oklahoma besides
those in the Oklahoma Aquarium in Jenks, OK. California has not had a
tornado above an EF2 over the last 10 years.

\subsection{Conclusions}\label{conclusions-1}

The probability of a Sharknado event is extremely low, if not a zero
percent chance. There was one tornado in the three coastal states that
might have had the ability to lift a shark. This South Carolina tornado
happened to start and end on land, making a Sharknado impossible. The
Oklahoma Aquarium houses sharks in a state that is known for its number
and strength of tornadoes. The Oklahoma Aquarium is in Tulsa County, OK,
where only EF2 or lower tornadoes have touched down in the last 10
years, rendering this very improbable. From the findings of this study,
the probability of a Sharknado happening is almost, if not, a zero
percent chance.

\subsection{Recommendations}\label{recommendations-1}

From this study, Sharknados are certainly extremely unlikely, if not a
zero percent chance. Even though strong storms and tornadoes without
sharks are still a cause for concern. Property damage and its resulting
insurance payouts will still be present with storms and tornadoes that
have high wind speeds and extreme precipitation events. These are real
and worrisome outcomes that happen all over the country that are more
probable than a Sharknado and should be taken more into account.




\end{document}
